\section{Passes}\label{sec:passes}

A \emph{pass} is a stretch of time during which every head moves monotonically, each in its own
direction. Transposition by recursive block exchange uses $O(\log n)$ passes (folklore), and the
external-memory intuition is that fewer passes cannot suffice. This section proves this for every
machine of the model whose running time is polynomial in $L$, whatever it does with its data. For
oblivious machines $r$ passes force $(k+1)^r\ge1+kn/(2+\sigma_M/b)$ (Theorem~P); for all machines
halting within $T$ steps they force $(k+1)^r\ge1+kn/(2+(\sigma_M+\gamma_T)/b)$ with
$\gamma_T=(k+4)\clog{(2T+1)}$, which is $O(kb)$ for $T\le\poly(L)$ (Theorem~P$'$). In both cases
$r\ge\log_{k+1}n-O(1)$ for $k\ge1$. Throughout,
\[
  \sigma_M=\clog{(q+1)}+k\,\clog{(g+1)}
\]
is the number of bits needed to write down a state and $k$ work symbols.

\subsection{Pass partitions and addresses}

\begin{definition}[Passes]\label{def:passes}
\formal{Tape/DefsP.lean}{MonoOn, AllMono, PassPartition, greedyBound, greedyPasses, addrCount}
A trajectory is \emph{monotone} on $[a,e)$ if all its steps in $[a,e)$ are non-decreasing, or all
are non-increasing. A \emph{pass partition} of the run of $M$ on $w$ into $r$ passes is a sequence
$0=p_0\le p_1\le\dots\le p_r=T$ such that on each $[p_i,p_{i+1})$ the input head and every work
head are monotone (empty passes are allowed). The \emph{greedy} partition extends each pass as far
as possible. We write
\[
  \Haddr=\Haddr(k,r)=\sum_{i<r}(k+1)^i ,\qquad\text{so that}\quad k\Haddr+1=(k+1)^r .
\]
\end{definition}

\begin{lemma}[Greedy passes]\label{lem:greedy}
\formal{Tape/Statements/Core.lean}{greedy_passPartition, greedyPasses_le, passPartition_iff, passPartition_indep, greedyPasses_indep, addrCount_mul}
The greedy passes form a pass partition, and a pass partition into $r$ passes exists if and only
if the greedy count is at most $r$. If $M$ is oblivious on $W$, the pass partitions of the runs on
the inputs of $W$ are the same for all inputs, and so is the greedy count. Moreover $k\Haddr+1=(k+1)^r$.
\end{lemma}

\begin{outline}
Minimality: the greedy partition stays ahead. By induction, the $i$-th greedy boundary is at least
the $i$-th boundary of any pass partition, since monotonicity on an interval is inherited by
subintervals. Conversely, a greedy partition with fewer passes is padded with empty passes. Independence
holds because pass partitions depend only on the head trajectories.
\end{outline}

\subsection{Where accessed cells come from}

\begin{lemma}[Fresh cells]\label{lem:fresh}
\formal{Tape/Statements/Core.lean}{pass_fresh}
Let $h$ be monotone on $[s,e)$ and $s\le\rho<a\le\tau\le e$. If $h(\tau)\ne h(a)$ then
$h(\rho)\ne h(\tau)$: during a monotone stretch, a cell occupied after time $a$, other than the
cell occupied at time $a$, was not occupied earlier in the stretch.
\end{lemma}

\begin{outline}
Monotonicity: $h(\rho)\le h(a)\le h(\tau)$ (or the reverse) with $h(a)\ne h(\tau)$.
\end{outline}

So, inside a pass, a work head that at time $\tau\ge a$ reads a cell other than the one it occupied at
time $a$ reads a cell last visited during $[a,\tau)$, a cell last visited in an \emph{earlier pass}, or a
blank cell. This is what makes the dependency structure of a pass schedule a forest of bounded size.

\subsection{The dependency forest}

Fix a run with a pass partition $p$ into $r$ passes and an output segment $[o_1,o_2)$.

\begin{definition}[Dependency forest]\label{def:forest}
\formal{Tape/DefsP.lean}{rootNode, sSet, childNode, forestNode, addrsBelow, addresses, forest, readSet, combSet}
\begin{itemize}[leftmargin=2em]
\item The \emph{root node} in pass $P_0$ is the set of steps of pass $P_0$ during which the output
length moves inside $[o_1,o_2)$.
\item For a node $J$, an earlier pass $P'$ and a work tape $t$, the \emph{S-set} consists of the
cells that head $t$ occupies during $J$, other than its cell at the first step of $J$, whose last
visit before that first step lies in pass $P'$. The \emph{child} of $J$ at $(P',t)$ is the set of
steps of pass $P'$ at which head $t$ lies between the least and the greatest cell of the S-set
(empty if the S-set is empty).
\item An \emph{address} is a pass $P_0<r$ followed by a list $(P_1,t_1),\dots,(P_j,t_j)$ with
$P_0>P_1>\dots>P_j$ and tapes $t_i$; the node at an address is obtained from the root in pass
$P_0$ by taking successive children. There are at most $\Haddr$ addresses
(\lean{card_addresses_le}). The \emph{forest} of the segment is the union of all its nodes.
\item $\mathrm{read}(J)$ is the set of input positions $<L$ read during $J$, and the \emph{comb}
of the segment is the set of input positions $d<n^2b$ with $\operatorname{tpos}(d)\in[o_1,o_2)$.
\end{itemize}
\end{definition}

Nodes are intervals of time, because heads are monotone within a pass and the output length is
monotone.

\begin{lemma}[Forest]\label{lem:forest}
\formal{Tape/Statements/Core.lean}{forest_subset, forest_entryPoints_subset, forest_entryCells_subset, forest_entryPoints_card_le, forest_entryCells_card_le, card_addresses_le, forest_info}
For every run with a pass partition into $r$ passes and every output segment, the forest consists
of steps $<T$; every entry point of the forest is the first step of a node; and every entry cell
of the forest is a read, at the first step of a node, of the cell under a work head. Hence
$\eone(\mathrm{forest})\le\Haddr$ and $\etwo(\mathrm{forest})\le k\Haddr$. Consequently, if $M$ is
oblivious on $\mathrm{Mat}(n,b)$ with halting time $T$ and transposes, then for the common run
\[
  2^{\abs{\mathrm{comb}\setminus\mathrm{read}(\mathrm{forest})}}\;\le\;
  \bigl((q+1)(g+1)^k\bigr)^{\#\mathrm{addresses}} .
\]
\end{lemma}

\begin{outline}
Entry points of a union of intervals are first steps of the intervals. For entry cells, take a
read by head $t$ at a step $\tau$ of a node $J$ that is not its first step $a$, of a cell whose last visit
before $\tau$ is at $\sigma$. If $\sigma\ge a$ then $\sigma\in J$, since $J$ is an interval. Otherwise the
cell is not the one occupied at time $a$ (that visit would come after $\sigma$), so by \cref{lem:fresh} it
was not visited earlier in the current pass: $\sigma$ lies in an earlier pass $P'$, the cell belongs to the
S-set at $(P',t)$, and $\sigma$ is a step of the child at $(P',t)$, which is part of the forest. So the read
is not external. For the last inequality, every output bit of the segment is
appended in some root node, so the comb bits that the forest does not read are in its cut set,
and the Cut Lemma on a set (\cref{lem:cut-on}) applies to the forest.
\end{outline}

\subsection{Disjointness and coverage}

\begin{lemma}[Disjointness]\label{lem:disjoint}
\formal{Tape/Statements/Core.lean}{forestNode_disjoint, read_chain}
At a fixed address, the nodes of two output segments, the first ending before the second starts,
are disjoint. If finitely many (at least one) segments have pairwise ordered output ranges (of any two,
one ends before the other starts), the input positions read by their nodes at a fixed address satisfy
$\sum_s\abs{\mathrm{read}(J_s)}+1\le L+\#\{s\}$.
\end{lemma}

\begin{outline}
By induction on the depth of the address, the nodes of an earlier and a later segment are
\emph{time-separated}, one lying entirely before the other: for the roots this is the monotonicity
of the output length. If two nodes of one pass are time-separated, then, the head being monotone in
that pass, every cell of the earlier node lies weakly on one side of the cell $c$ under the head at the
later node's start, while the S-set of the later node, which excludes $c$, lies strictly on the other side;
so the hulls of the two S-sets are disjoint, and monotonicity in the earlier pass makes the children
time-separated. For the second part, all nodes at a fixed address lie in one pass, in which the input head
is monotone; ordered in time, the read set of each node meets the union of the read sets of the earlier
nodes in at most one position (the input cell at its first step), so
$\sum_s\abs{\mathrm{read}(J_s)}\le\abs{\bigcup_s\mathrm{read}(J_s)}+\#\{s\}-1\le L+\#\{s\}-1$.
\end{outline}

Group the columns into $\lceil n/\gamma\rceil$ groups of $\gamma$ consecutive columns (the last may be
smaller); group $s$ owns the output segment of its columns, and its comb consists of at most $\gamma b$
consecutive input bits in each of the $n$ rows.

\begin{lemma}[Coverage]\label{lem:coverage}
\formal{Tape/Statements/Core.lean}{coverage_per_address, coverage_grouped}
For every run with a pass partition into $r$ passes and every $\gamma\ge1$, at a fixed address the comb
bits read by the nodes of the $\lceil n/\gamma\rceil$ groups number at most $(\gamma+1)nb$ in total.
Over all addresses, the comb bits read by the group forests number at most
$\#\mathrm{addresses}\cdot(\gamma+1)nb$.
\end{lemma}

\begin{outline}
At a fixed address the read sets of the groups form a time-ordered chain of intervals of input
positions (\cref{lem:disjoint}). An exact residue count bounds what such a chain can collect from
the combs: with $m=nb$, the positions $d$ of a set $R$ of consecutive positions with $d\bmod m$ in
a set $\mathrm{Res}$ number at most $\abs{\mathrm{Res}}(\abs{R}+m)/m$; summing over the chain with
the true group widths gives $(\gamma+1)nb$.
\end{outline}

\subsection{Theorem P}

\begin{theorem}[Theorem P$_\gamma$, grouped columns]\label{thm:Pgamma}
\formal{Tape/Statements/Core.lean}{theoremP_grouped, theoremP_grouped_bits}
Let $M$ be oblivious on $\mathrm{Mat}(n,b)$ with halting time $T$, transpose $n\times n$ matrices
with $b$-bit entries, and let the common run have a pass partition into $r$ passes. Then for every
$\gamma\ge1$
\[
  2^{n^2b}\le 2^{\Haddr(\gamma+1)nb}\bigl((q+1)(g+1)^k\bigr)^{\lceil n/\gamma\rceil\Haddr},
  \qquad
  n^2b\le\Haddr\bigl((\gamma+1)nb+\lceil n/\gamma\rceil\,\sigma_M\bigr).
\]
\end{theorem}

\begin{theorem}[Theorem P]\label{thm:P}
\formal{Tape/Statements/Core.lean}{theoremP, theoremP_bits, theoremP_bits_add_one, theoremP_div, theoremP_pow, theoremP_greedy}
Under the hypotheses of \cref{thm:Pgamma},
\[
  2^{n^2b}\le 2^{2nb\Haddr}\bigl((q+1)(g+1)^k\bigr)^{n\Haddr},\qquad
  nb\le\Haddr\,(2b+\sigma_M),\qquad
  k\,nb\le\bigl((k+1)^r-1\bigr)(2b+\sigma_M),
\]
also in the weaker form $n^2b\le\Haddr(n(2b+\sigma_M)+1)$, and the same holds with $r$ the greedy pass
count (the greedy passes form a pass partition,
\cref{lem:greedy}; \lean{theoremP_greedy} states the bits form). In particular, for $k\ge1$ and
$b=\lceil\log_2n\rceil$, a fixed machine needs $r\ge\log_{k+1}n-O(1)$ passes.
\end{theorem}

\begin{outline}
Theorem~P is the case $\gamma=1$ of Theorem~P$_\gamma$. For each group $s$, \cref{lem:forest}
bounds the comb bits of $s$ not read by the group's forest:
$2^{\abs{\mathrm{comb}_s\setminus U_s}}\le((q+1)(g+1)^k)^{\Haddr}$. \Cref{lem:coverage} bounds
the comb bits that are read: $\sum_s\abs{\mathrm{comb}_s\cap U_s}\le\Haddr(\gamma+1)nb$. The combs
partition the $n^2b$ input positions, so
$n^2b\le\sum_s\abs{\mathrm{comb}_s\setminus U_s}+\Haddr(\gamma+1)nb$. The per-group bounds are summed,
$\sum_s2^{\abs{\mathrm{comb}_s\setminus U_s}}\le\lceil n/\gamma\rceil((q+1)(g+1)^k)^{\Haddr}$, and the
inequality of arithmetic and geometric means (the final step shared with Theorem~P$'$) gives
$2^{\sum_s\abs{\mathrm{comb}_s\setminus U_s}}\le((q+1)(g+1)^k)^{\lceil n/\gamma\rceil\Haddr}$, hence the
claim. The other forms are arithmetic, using $k\Haddr+1=(k+1)^r$.
\end{outline}

Theorem~P allows arbitrary coding: the machine may combine bits in its finite control and on its
tapes in any way, dwell, and move its heads at arbitrary relative speeds. The closest previous lower
bound for transposition is Kirchherr's $\Omega(\log n)$ reversals for \emph{conservative} machines, which
only move bits~\cite{Kirchherr1988}; the stream lower bounds~\cite{GroheSchweikardt2005,GHS2009} need long or
index-tagged items (\cref{sec:intro}).

\subsection{Machines correct on part of the input}

\begin{theorem}[Theorem P on a set]\label{thm:P-on}
\formal{Tape/Statements/Core.lean}{theoremP_on, theoremP_on_bits, theoremP_on_div, theoremP_on_pow, theoremP_grouped_on, theoremP_grouped_on_bits, theoremP_on_robust_gamma, theoremP_on_robust}
Let $M$ be oblivious on a set $W$ with halting time $T$ and correct on $W$, let the common run
have a pass partition into $r$ passes, and let $W$ have co-entropy at most $m$. Then for every
$\gamma\ge1$
\[
  n^2b\le\lceil n/\gamma\rceil\,m+\Haddr\bigl((\gamma+1)nb+\lceil n/\gamma\rceil\,\sigma_M\bigr):
\]
each group of $\gamma$ columns pays $m$ bits. In particular:
\begin{enumerate}[label=(\alph*)]
\item at co-entropy $Cn$ ($C\in\N$): $nb\le Cn+\Haddr(2b+\sigma_M)$, so $\Haddr\ge n(b-C)/(2b+\sigma_M)$, and
$k\,nb\le kCn+((k+1)^r-1)(2b+\sigma_M)$;
\item at co-entropy $Cnb$ ($C\in\N$), for every $\gamma\ge1$:
$(\gamma-C)\,n^2b\le\gamma C\,nb+\Haddr\bigl(\gamma(\gamma+1)nb+(n+\gamma)\sigma_M\bigr)$; with
$\gamma=2C$, $Cnb\le 2C^2b+\Haddr\bigl(2C(2C+1)b+2\sigma_M\bigr)$, so
$\Haddr\gtrsim n/(2(2C+1))$ and, for $k\ge1$, still $r\ge\log_{k+1}n-O(\log C)$.
\end{enumerate}
\end{theorem}

\begin{outline}
As for Theorem~P$_\gamma$, with the co-entropy term of the Cut Lemma on a set (\cref{lem:cut-on}) kept in
each group's bound: each group's forest pays the co-entropy $m$ once. In the formal development this robust
form is the only direct proof; Theorem~P$_\gamma$ is its case $W=\mathrm{Mat}(n,b)$, $m=0$. The
corollaries are arithmetic.
\end{outline}

\subsection{Data-dependent heads: Theorem P\texorpdfstring{$'$}{'}}

For a machine whose head motion depends on the input, every input has its own run, pass partition
and dependency forests. A forest is then described by a \emph{structure code} per node: its first
and last steps, the positions of all heads and the output length at its first step
(\lean{nodeCode}); and by a \emph{side code}: the state and the symbols under the work heads at its
first step (\lean{sideCode}). The structure codes of nodes inside $[0,T)$ take at most
$(2T+1)^{k+4}$ values (\lean{structCount} is this bound; the count is part of the proof of
\cref{lem:forest-agree}).

\begin{lemma}[Forests of different inputs]\label{lem:forest-agree}
\formal{Tape/Statements/Adaptive.lean}{locality_anchored, forest_agree, forest_comb_agree, kraft_group}
Let $w,w'$ be inputs of the same length $L$, with pass partitions of $[0,T)$ and of $[0,T')$ into $r$
passes (each its own), and suppose their forests of a segment have equal structure and side codes at every
address, and that $w$ and $w'$ agree on every input position $<L$ read by both forests. Then the two
forests have the same nodes, and the runs agree on them (state, positions, output length, symbols under the
heads, transitions). If $M$ transposes $n\times n$ matrices with $b$-bit entries, $L=n^2b$, and $w,w'$ halt
at $T$ and $T'$ respectively, they agree on the segment's comb. Consequently, if $M$ transposes and every
$v\in\{0,1\}^{n^2b}$ halts at some $T_v\le T$ with a pass partition of $[0,T_v)$ into $r$ passes, then for
every segment
\[
  \sum_{v\in\{0,1\}^{n^2b}}2^{\abs{\mathrm{comb}\setminus\mathrm{read}_v(\mathrm{forest}_v)}}
  \;\le\;2^{n^2b}\Bigl((2T+1)^{k+4}(q+1)(g+1)^k\Bigr)^{\Haddr}.
\]
\end{lemma}

\begin{outline}
Nodes are intervals and the structure code records their first and last steps, so equal codes give equal
nodes at every address: the two forests are the same set of steps. The first part is then an anchored form
of the Locality Lemma (\lean{locality_anchored}, a corollary of \cref{lem:locality-core}) with the first
steps of the nodes as anchors: the structure and side codes give agreement at the anchors, and
\cref{lem:forest}, applied to each run, shows that every other step of the forest has its predecessor in
the forest and that every cell read there was last written inside the forest (or never); the input symbols
agree at commonly read positions because the inputs have the same length. The comb then agrees by
correctness. For
the Kraft-type inequality, fix the codes: within a class of inputs with equal codes, overwriting
the unread comb bits of $v$ by arbitrary values is injective, because two overwritten inputs that
coincide agree off the union of their free positions, hence on every position read by both
forests, hence (by the first part) on the comb, hence everywhere.
\end{outline}

\begin{theorem}[Theorem P$'$]\label{thm:Pprime}
\formal{Tape/Statements/Adaptive.lean}{theoremPprime_grouped, theoremPprime_grouped_bits, theoremPprime, theoremPprime_bits, theoremPprime_div, theoremPprime_pow, theoremPprime_greedy, theoremPprime_gamma, theoremPprime_ratio}
Let $M$ transpose $n\times n$ matrices with $b$-bit entries, and suppose every matrix input halts
within $T$ steps with a pass partition of its own run into $r$ passes. Let
$\gamma_T=(k+4)\clog{(2T+1)}$. Then for every $\gamma\ge1$
\[
  n^2b\le\Haddr\bigl((\gamma+1)nb+\lceil n/\gamma\rceil(\sigma_M+\gamma_T)\bigr),
  \qquad\text{in particular}\qquad nb\le\Haddr\,(2b+\sigma_M+\gamma_T),
\]
and $k\,nb\le((k+1)^r-1)(2b+\sigma_M+\gamma_T)$; these hold also in log-free form and with greedy pass
counts. The bound also holds with
$\gamma'=\clog r+2(k+2)\clog{(2T+3)}\ge\gamma_T$ in place of $\gamma_T$, in the forms
$n^2b\le\Haddr(n(2b+\sigma_M+\gamma')+1)$ and, for $b\ge1$,
$\Haddr\ge n/(2+(\sigma_M+\gamma'+1)/b)$.
\end{theorem}

\begin{outline}
By \cref{lem:coverage}, which holds run by run,
$\sum_s\abs{\mathrm{comb}_s\setminus U_s(v)}\ge n^2b-\Haddr(\gamma+1)nb$ for every input $v$. Sum
the inequality of \cref{lem:forest-agree} over the groups and average over the $2^{n^2b}$ inputs:
some $v_0$ has $\sum_s2^{m_s(v_0)}\le\lceil n/\gamma\rceil K$, where $m_s(v)$ is the number of
unread comb bits of group $s$ and $K=((2T+1)^{k+4}(q+1)(g+1)^k)^{\Haddr}$. The inequality of
arithmetic and geometric means gives $2^{\sum_sm_s(v_0)}\le K^{\lceil n/\gamma\rceil}$. The
remaining forms are arithmetic.
\end{outline}

If $T\le\poly(L)$ and $b=\lceil\log_2n\rceil$, then $\gamma_T=O(kb)$ and Theorem~P$'$ again gives
$r\ge\log_{k+1}n-O(1)$; this covers every machine allowed by~(T). Theorem~P$'$ also has the robust
form of \cref{thm:P-on}.

\begin{theorem}[Theorem P$'$ on a set, and against the frozen statement]\label{thm:Pprime-on}
\formal{Tape/Statements/Adaptive.lean}{forest_comb_agree_on, kraft_group_on, theoremPprime_grouped_on, theoremPprime_grouped_on_bits, theoremPprime_on, theoremPprime_on_bits, theoremPprime_on_div, theoremPprime_on_pow, theoremPprime_on_robust, theoremPprime_on_robust_gamma, theoremPprime_on_greedy, theoremPprime_linearTime}
Suppose that on every input of a set $W$ of co-entropy at most $m$, the machine $M$ has, at some
time $T'\le T$, the transposed output together with a pass partition of $[0,T')$ into $r$ passes
(it need not halt; other inputs are unconstrained). Then
$n^2b\le\lceil n/\gamma\rceil m+\Haddr((\gamma+1)nb+\lceil n/\gamma\rceil(\sigma_M+\gamma_T))$ for
every $\gamma\ge1$, and the inequalities of \cref{thm:P-on}(a),(b) hold with $\sigma_M+\gamma_T$ in place
of $\sigma_M$. Their asymptotic consequences change: $T\ge L$ forces $\gamma_T\ge2(k+4)\log_2n$, which is
not $o(b)$; if $\gamma_T\le c'b$ (for instance $T\le\poly(L)$), (b) gives
$\Haddr\ge(Cn-2C^2)b/(2C(2C+1)b+2(\sigma_M+\gamma_T))$ and, for $k\ge1$, $r\ge\log_{k+1}n-O(\log(C+c'))$. If $M$ transposes within $c(L+1)$ steps (\cref{def:T}) and every run on an
$n\times n$ input has greedy pass count at most $r$, then
$n^2b(n)\le\Haddr\,n\bigl(2b(n)+\sigma_M+(k+4)\clog{(2c(L+1)+1)}\bigr)$.
\end{theorem}

\begin{outline}
As for Theorem~P$'$, with the Kraft sum over $W$: the injectivity step uses correctness only on
the two inputs it compares, and the averaging step uses $\abs{W}\ge2^{n^2b-m}$, which charges $m$
once per group. The last statement instantiates $T=c(L+1)$.
\end{outline}

\subsection{Non-vacuity and non-triviality}

\begin{proposition}\label{prop:P-witness}
\formal{Tape/Statements/Core.lean}{theoremP_hyps_tr2, theoremP_hyps_tr2k}
The hypotheses of Theorem~P are satisfiable, with $k=0$ (\lean{Witness.tr2}, $n=2$, $b=1$) and
with $k=1$ (\lean{Witness.tr2k}).
\end{proposition}

\begin{outline}
Evaluation of the concrete machines on all inputs; the greedy passes give the pass partition
(\cref{lem:greedy}).
\end{outline}

\begin{proposition}\label{prop:Pprime-witness}
\formal{Tape/Statements/Adaptive.lean}{theoremPprime_hyps_of_oblivious, theoremPprime_hyps_tr2dd, theoremPprime_no_five_pass}
\begin{enumerate}[label=(\alph*)]
\item Every machine that is oblivious on $\mathrm{Mat}(n,b)$ with halting time $T$ and whose common run
has a pass partition into $r$ passes satisfies the run hypothesis of Theorem~P$'$ (every matrix input halts
within $T$ steps with a pass partition into $r$ passes). In particular the machines of Theorem~P satisfy all
hypotheses of Theorem~P$'$, with the same $T$ and $r$.
\item A machine \lean{Witness.tr2dd} transposes $2\times2$ matrices with $1$-bit entries, is
oblivious for no halting time (its running time is $5$ or $6$ depending on the input), and every
input halts within $6$ steps in a single pass.
\item No machine with one work tape, at most $2^8$ states and at most $4$ work symbols transposes
$1024\times1024$ matrices with $10$-bit entries so that every input halts within $2^{30}$ steps in
$5$ passes.
\end{enumerate}
\end{proposition}

\begin{outline}
(a) by \cref{lem:greedy}; (b) by evaluation; (c) by Theorem~P$'$ and arithmetic.
\end{outline}
